\section{Cell energy and a principal-to-diagonal transfer}
\label{sec:diagonal-transfer}

TPC-100 first returns the \(q_X\mid u\) subcarrier at
\(O(X^{o(1)}Q_X^2)\).  It then refines the remaining
\(q_X\nmid u\) mass into opened atoms and groups those atoms into
lossless cells while retaining the parent branch length
\(N_\beta\).  Write \(a_H^{\mathrm{inv}}\) for this invertible-ultra
census.  For every exact width \(H\), TPC-100 proves
\begin{equation}
 \|(a_H^{\mathrm{inv}})^\circ\|_2^2
 =
 \mathfrak D_H^\circ+\mathfrak X_H^\circ,
\qquad
 \mathfrak D_H^\circ\ge0,
\label{eq:tpc100-energy}
\end{equation}
where \(\mathfrak D_H^\circ\) is the one-cell diagonal ledger and
\(\mathfrak X_H^\circ\) is the centered excess from collisions
between distinct cells.

At the maximally literal TPC-100 resolution, a cell retains
\((m,u,\sigma)\), so the source identity recovers
\[
 j=\frac{\sigma u-h_0}{m}.
\]
Every such fine cell is therefore a singleton.  Consequently
\[
 \mathfrak D_H^\circ
 =
 \frac{q_X-2}{q_X-1}
 \sum_{\substack{p=(\beta,z):\,H_\beta=H\\
                  q_X\nmid u_p}}
 w_{\beta,z}^{\,2}.
\]
This is an exact atomic diagonal ledger, not an intra-cell
dispersion gain; every collision between distinct opened atoms
remains in \(\mathfrak X_H^\circ\).

\begin{theorem}[Width-weighted cell gates]
\label{thm:width-cell-gates}
After the \(q_X\mid u\) return, the full principal condition
\eqref{eq:new-principal-gate} (or its restriction to the invertible
census), together with the two estimates
\begin{align}
 \sum_Hg_{q_X}(H)\sqrt{\mathfrak D_H^\circ}
 &\ll X^{o(1)}Q_X^2,
\label{eq:weighted-diagonal-gate}\\
 \sum_Hg_{q_X}(H)\sqrt{|\mathfrak X_H^\circ|}
 &\ll X^{o(1)}Q_X^2
\label{eq:weighted-cross-gate}
\end{align}
are sufficient to return the remaining invertible positive
nonconstant resonance carrier.  Adding the already returned
\(q_X\mid u\) piece gives
\(O(X^{o(1)}Q_X^2)\) for the complete positive nonconstant carrier.
\end{theorem}

\begin{proof}
By \eqref{eq:tpc100-energy},
\[
 \|(a_H^{\mathrm{inv}})^\circ\|_2
 \le
 \sqrt{\mathfrak D_H^\circ}
 +
 \sqrt{|\mathfrak X_H^\circ|}.
\]
Multiply by \(g_{q_X}(H)\), sum over \(H\), and apply
\cref{thm:literal-width-bound}.
\end{proof}

There is an abstract way to discharge the diagonal condition from
the principal ledger.  It isolates exactly the additional local
fact that must be checked in the literal weights.

\begin{definition}[Opened-atom cap]
\label{def:atom-cap}
Let \(w_C(j)\ge0\) be the normalized \(q_X\nmid u\) opened-atom
weights used in the TPC-100 cell census, so
\[
 M_H^{\mathrm{inv}}
 =
 \|a_H^{\mathrm{inv}}\|_1
 =
 \sum_{C,j}w_C(j).
\]
Put
\[
 W_X=\max_{H,C,j}w_C(j)
\]
over the actual support, with the convention \(W_X=0\) when that
support is empty, and let
\(\mathcal H_X
=\{H:g_{q_X}(H)>0,\ M_H^{\mathrm{inv}}>0\}\).
\end{definition}

\begin{theorem}[Principal-to-diagonal transfer]
\label{thm:principal-diagonal-transfer}
Define the principal mass without the bounded filter factor by
\[
 \mathfrak P_X^{\mathrm{inv}}
 =
 \sum_H\frac{2H}{q_X-1}M_H^{\mathrm{inv}}.
\]
Then
\begin{equation}
 \boxed{
 \sum_Hg_{q_X}(H)\sqrt{\mathfrak D_H^\circ}
 \le
 \sqrt{
 W_X\,|\mathcal H_X|\,(q_X-1)\,
 \mathfrak P_X^{\mathrm{inv}}}.}
\label{eq:principal-diagonal-transfer}
\end{equation}
Consequently, if
\[
 q_X\asymp Q_X,\qquad
 W_X\le X^{o(1)},\qquad
 \mathfrak P_X^{\mathrm{inv}}\le X^{o(1)}Q_X^2,
\]
then \eqref{eq:weighted-diagonal-gate} follows automatically.
\end{theorem}

\begin{proof}
The diagonal definition and nonnegativity give
\[
 0\le\mathfrak D_H^\circ
 \le\sum_{C,j}w_C(j)^2
 \le W_XM_H^{\mathrm{inv}}.
\]
Therefore Cauchy--Schwarz over the active widths yields
\begin{align*}
 \sum_Hg_{q_X}(H)\sqrt{\mathfrak D_H^\circ}
 &\le
 \sqrt{W_X}
 \sum_{H\in\mathcal H_X}
 g_{q_X}(H)\sqrt{M_H^{\mathrm{inv}}}\\
 &\le
 \sqrt{
 W_X|\mathcal H_X|
 \sum_Hg_{q_X}(H)^2M_H^{\mathrm{inv}}}.
\end{align*}
Since
\[
 g_{q_X}(H)^2\le2H,
\]
the last sum is at most
 \((q_X-1)\mathfrak P_X^{\mathrm{inv}}\), proving
\eqref{eq:principal-diagonal-transfer}.  Finally
\(|\mathcal H_X|\le(q_X-1)/2\) and \(q_X\asymp Q_X\).
The full principal mass dominates
\(\mathfrak P_X^{\mathrm{inv}}\) because the opened weights are
nonnegative.
\end{proof}

\begin{remark}[A real reduction, not a hidden proof]
Once the literal atom cap \(W_X\le X^{o(1)}\) and the principal gate
are verified, no separate diagonal-energy theorem is needed.  This
paper does not assert that cap: an opened atom includes its actual
row coefficients, exact-content and projector factors, smooth
weights, and the single global normalization.  Bounding that
specific quantity is the next literal audit, not a formal
consequence of row \(\ell^1\) mass.
\end{remark}

\begin{remark}[The cross-cell door remains]
The transfer does not control
\(\mathfrak X_H^\circ\).  TPC-100 shows that its support obeys an
exact affine collision congruence, but distinct cells can still
align maximally.  Since the fine cells are singleton, this includes
all collisions between distinct atoms, even two atoms under the
same parent branch.  The surviving positive route is therefore the
weighted cross-cell gate \eqref{eq:weighted-cross-gate}, or a return
to signed arithmetic before the positive census is formed.
\end{remark}
